\documentclass[11pt]{article}
\usepackage[T1]{fontenc}
\usepackage{lmodern,amsmath,amssymb,amsthm,booktabs,tabularx}
\usepackage[margin=1in]{geometry}
\usepackage{xurl}
\usepackage[hidelinks]{hyperref}
\newtheorem{proposition}{Proposition}[section]
\newtheorem{corollary}[proposition]{Corollary}
\theoremstyle{remark}
\newtheorem{remark}[proposition]{Remark}
\newcommand{\Ric}{\operatorname{Ric}}
\newcommand{\Vol}{\operatorname{Vol}}
\newcommand{\CP}{\mathbb{CP}}
\newcommand{\E}{\mathcal E}
\newcommand{\J}{\mathcal J}
\newcommand{\avint}{\frac{1}{V}\int_M}
\title{Classical Curvature Bounds and Quantitative Einstein Gaps\\
\large Inheritance, differential coercivity, and extremal constructions}
\author{}
\date{4 October 2026}
\begin{document}
\maketitle
\begin{abstract}
Classical monopole estimates already imply normalized Einstein gaps away
from a critical topological ratio, and classical surgery constructions
already provide matching bounds for certain unnormalized curvature
functionals. We derive these implications in a common normalization and
identify the additional content of uniform critical separation and linear
coercivity. Explicit formal integral data distinguish the latter from the
listed classical numerical inequalities. We also record the subsequent
finite-scale construction that lowers the candidate balanced coefficient
below four, without identifying it with the unrestricted optimum.
\end{abstract}

\section{The common problem and the objects being compared}

The problem is how closely a smooth metric can satisfy the Einstein
equation under specified topological and actual-field constraints.
Topology supplies integral relations. An actual field also supplies
differential relations with the curvature of the \emph{same} metric.
The question is whether these coupled relations impose quantitative
constraints beyond the numerical curvature budgets, and which geometries
realize the resulting bounds. Obstructions, constructions, variation,
degeneration and exact realization address different parts of this question.

On a closed oriented smooth four-manifold, write
\begin{gather}
 V=\Vol(M,g),\quad I=\int_Ms\,dV_g,\quad
 S=\int_Ms^2\,dV_g,\quad P=\int_M|\Ric^\circ|^2\,dV_g,\\
 R=\int_M|\Ric|^2\,dV_g=P+S/4,\qquad
 W_\pm=\int_M|W^\pm|^2\,dV_g,\qquad T_\pm=2\chi\pm3\tau.
\end{gather}
Weyl norms are Hilbert--Schmidt norms of operators on two-forms, not
four-index tensor norms. Thus
\begin{equation}\label{eq:cw}
 4\pi^2T_\pm=S/24+2W_\pm-P/2,
 \qquad \int_M|\mathrm{Rm}_{\mathrm{op}}|^2=8\pi^2\chi+P.
\end{equation}
Set
\begin{equation}
 \E(g)=\inf_{\lambda\ne0}
 \frac{\int_M|\Ric-\lambda g|^2}{\lambda^2V},
 \qquad \J(g)=\frac{4P}{R}\quad(R>0).
\end{equation}
For $R>0$, completing the square in $1/\lambda$ gives
\begin{equation}\label{eq:defects}
 \E=4-\frac{I^2}{VR}
     =\J+\frac{S-I^2/V}{R}\ge\J.
\end{equation}
This remains an infimum identity when $I=0$, although no finite
$\lambda$ minimizes the expression. Constant scalar curvature gives
$\E=\J$. For a Ricci-flat metric, $R=0$, $\J$ is undefined and
$\E=4$; exclusions of Ricci-flat metrics in positive-budget gauge
applications therefore use the actual equations separately.

In particular, $\inf R$, $\inf P$ and $\inf\J$ are different variational
problems. The pointwise identity $P(g)=R(g)-S(g)/4$ does not imply
\begin{equation}\label{eq:not-inf}
 \inf_gP(g)=\inf_gR(g)-\tfrac14\inf_gS(g).
\end{equation}
A common recovery sequence gives an upper bound on $\inf P$, not a
lower bound against other competitors.

\section{Normalized gaps already implied by classical estimates}

Let $a$ be an actual determinant class with
$p=(a_g^+)^2\ge A>0$. In the monopole or conformal-solvability range of
the classical estimates \cite{LeBrunRicci}, one has
\begin{equation}\label{eq:classical}
 S\ge32\pi^2p,\qquad
 \sqrt S+\sqrt{6W_+}\ge6\sqrt2\,\pi\sqrt p.
\end{equation}
These bounds do not require a diameter, noncollapse or uniform Sobolev
assumption. They already use actual Seiberg--Witten equations and
Weitzenb\"ock identities.

Weighted Cauchy--Schwarz yields, for every $\beta>0$,
\begin{equation}\label{eq:Dbeta}
 D_\beta=(1+\beta^{-1})S+6(\beta+1)W_+-72\pi^2p\ge0.
\end{equation}
With $T=T_+$, substitution of \eqref{eq:cw} gives exactly
\begin{equation}\label{eq:Dcw}
 D_\beta=(\beta+1)
 \left[\frac{8-\beta}{8\beta}S+\frac32P+12\pi^2T\right]
 -72\pi^2p.
\end{equation}

\begin{proposition}[A classical subcritical ratio bound]
If $0<t:=T/A<2/3$ and \eqref{eq:classical} holds, then
\begin{equation}\label{eq:old-gap}
 \J(g)\ge\frac{2-3t}{2-t}.
\end{equation}
If $T\le0$ and $R>0$, then $\J(g)\ge1$ directly from
\eqref{eq:cw}.
\end{proposition}
\begin{proof}
Choose $\beta+1=6/t$, so $\beta>8$ and
$12\pi^2(\beta+1)T=72\pi^2A\le72\pi^2p$. Consequently
\[
 0\le D_\beta\le(\beta+1)
 \left[\frac{8-\beta}{8\beta}S+\frac32P\right],
 \qquad \frac PS\ge\frac{\beta-8}{12\beta}.
\]
Here $S>0$ by the scalar estimate. Substituting in
$\J=4(P/S)/(P/S+1/4)$ proves \eqref{eq:old-gap}.
For $T\le0$, \eqref{eq:cw} gives $P\ge S/12+4W_+$, hence
$R\le4P$.
\end{proof}
Thus, for $0<\delta\le2/3$,
\begin{equation}
 T/A\le2/3-\delta\quad\Longrightarrow\quad
 \J\ge\frac{9\delta}{4+3\delta}.
\end{equation}
This is an algebraic consequence of classical joint estimates, not a new
analytic theorem. A lower bound for $P$ alone cannot simply be divided by
an uncontrolled $R$, but the \emph{joint} classical estimates do give a
normalized ratio bound.

\section{The critical separation and its precise input}

At $t=2/3$, the preceding lower bound vanishes. The aligned parameter is
$\beta=8$, for which
\begin{equation}\label{eq:D8}
 D_8=108\pi^2(T-2p/3)+\frac{27}{2}P.
\end{equation}
When $p=A$ and $T=2A/3$, the numerical statement $D_8\ge0$ reduces to
$P\ge0$.

Classical equality analysis is stronger than this numerical statement.
The convex-hull theory identifies the relevant equality geometries
\cite{LeBrunConvex}; the strict inequality in Theorem 7.5 of the 2026
overview \cite{LeBrunOverview} excludes equality in its curvature bound
when its convex budget is positive. That theorem does not state a uniform
positive improvement. The distinction is
\[
 (\forall g)\ \J(g)>0\quad\not\Longrightarrow\quad
 \inf_g\J(g)>0.
\]
The 2006 convex-geometry paper and the 2026 overview are different
references; Theorem 7.5 belongs to the latter.

\subsection*{The quantitative input supplied by the current manuscript}
On a negative Yamabe representative $\widehat g$ with $s_{\widehat g}=-4$,
the single-solution parameter theorem supplies
\begin{equation}\label{eq:raw}
 (1+b)\left(r+\frac{D_8}{8V}\right)\ge c_{8,\mathrm{raw}}>0,
 \quad r=P(\widehat g)/V,\quad b=\frac{(-a^2)_+}{V}.
\end{equation}
Its input is the actual unperturbed monopole solution for the specified
structure, with $p\ge A>0$; the same structure is retained throughout
the argument. This analytic theorem, not \eqref{eq:Dbeta} alone, is the
additional ingredient. The full input and proof are in
\path{research/paper/sections/tracefree_gap.tex} and
\path{research/paper/sections/single_solution_coercivity.tex}, relative
to the Skill root.

For completeness, the algebraic transfer is short. Put
$d=(T/A-2/3)_+$ and $\nu=(-a^2)_+/A$. The scalar budget gives
$A/V\le1/(2\pi^2)$, and \eqref{eq:D8} implies
\begin{equation}\label{eq:raw-upper}
 r+\frac{D_8}{8V}\le\frac{43}{16}r+\frac{27}{4}d.
\end{equation}
For $\delta=S(g)-S(\widehat g)\ge0$, conformal invariance of Weyl
energy and \eqref{eq:cw} give
\begin{equation}\label{eq:conformal}
 P(g)=P(\widehat g)+\delta/12,\quad
 R(g)=R(\widehat g)+\delta/3,
\end{equation}
\[
 \J(g)=\frac{R(\widehat g)}{R(\widehat g)+\delta/3}\J(\widehat g)
        +\frac{\delta/3}{R(\widehat g)+\delta/3}.
\]
If $\J(g)<1$, this yields $\J(\widehat g)\le\J(g)<1$ and
$r=4\J(\widehat g)/(4-\J(\widehat g))\le4\J(g)/3$.
Since $b\le\nu/(2\pi^2)$, \eqref{eq:raw}--\eqref{eq:raw-upper}
give the manuscript's N1 conclusion
\begin{equation}\label{eq:N1}
 \boxed{\J(g)+d\ge\frac{c_*}{1+\nu/(2\pi^2)},\qquad
 c_*:=\min\{1,4c_{8,\mathrm{raw}}/27\}.}
\end{equation}
The branch $\J(g)\ge1$ is immediate. In particular $a^2\ge0$
already removes this negative-square cost; requiring $a^2\ge T$ is
sufficient in the previously stated version but is not necessary for
\eqref{eq:N1}.

The relative-shortfall version is also available. With
$y=(T-a^2)_+/A$, \eqref{eq:cw} yields
\[
 \frac{(-T)_+}{V}\le\frac r{8\pi^2},\qquad
 b\le\frac r{8\pi^2}+\frac y{2\pi^2}.
\]
In the $\J<1$ branch $r<4/3$, so the same proof gives
\begin{equation}\label{eq:N1-y}
 \J(g)+d\ge\frac{c}{1+y}.
\end{equation}
The positive gap for $\J$ itself is the right-hand side minus $d$.
For example a uniform $\varepsilon$ in the zero-cost branch gives
$\J\ge\varepsilon/2$ when $T/A\le2/3+\varepsilon/2$.
Thus the additional conclusion concerns a neighborhood of the critical
ratio, not only equality at one value.

An all-metric application must supply the stated solution on every
negative Yamabe representative with uniform budgets. Older conformal
supply and positive projection imply Yamabe negativity by the ordinary
monopole Weitzenb\"ock identity: a nonnegative constant-scalar
representative forces the spinor to vanish, then $p=0$.
Pin-minus descent uses actual lifted solutions and lifted metrics; it
does not require arbitrary non-invariant solutions on the cover.

\section{Classical recovery and what it does not minimize}

For a minimal complex surface of general type $X$, put
\[
 M=X\#k\overline{\CP}^{\,2}\#\ell(S^1\times S^3),\qquad
 A=c_1^2(X),\quad Q=k+4\ell,\quad h=Q/A.
\]
LeBrun's Theorem 4.3 \cite{LeBrunRicci} gives
\begin{equation}
 \inf_gR(g)=8\pi^2(A+Q).
\end{equation}
The recovery construction simultaneously has
\[
 S\longrightarrow32\pi^2A,\qquad
 P\longrightarrow8\pi^2Q,
\]
and vanishing scalar variance in the corresponding normalized recovery.
Consequently
\begin{equation}\label{eq:old-upper}
 \inf_g\J(g)\le\frac{4h}{1+h},\qquad
 \inf_g\E(g)\le\frac{4h}{1+h}.
\end{equation}
These upper bounds inherit classical constructions. Related BF
connected-sum results also compute Ricci-energy infima by matching
lower estimates and constructions \cite{IshidaLeBrun}.

For the handle-free case, the full-curvature bounds in Section 11 of
\cite{LeBrunOptimal}, combined with \eqref{eq:cw}, imply
\begin{equation}
 8\pi^2\max\{0,k-A/3\}\le\inf_gP(g)\le8\pi^2k.
\end{equation}
Neither this interval nor the exact value of $\inf R$ determines
$\inf P$. Equation \eqref{eq:not-inf} cannot be used to close that gap.

\begin{proposition}[Equality of the two low-defect infima]
Suppose all conformal classes under consideration are negative Yamabe
classes and the class of admissible metrics is closed under their
Yamabe replacements. If $\inf_g\E(g)<1$, then
$\inf_g\E(g)=\inf_g\J(g)$.
\end{proposition}
\begin{proof}
Equation \eqref{eq:defects} gives $\inf\J\le\inf\E<1$.
Take a $\J$-minimizing sequence eventually below one. The convex
combination in \eqref{eq:conformal} gives
$\J(\widehat g_i)\le\J(g_i)$. Since
$\E(\widehat g_i)=\J(\widehat g_i)$, the reverse inequality of
infima follows.
\end{proof}
In the inherited small-$h$ family, \eqref{eq:old-upper} makes this
applicable for $h<1/3$. This is not an unrestricted identity of infima
for all manifolds or energy ranges.

\section{Why the linear constraint contains additional information}

The manuscript's I40/I41 chain and its single-solution continuation give
\begin{equation}\label{eq:linear}
 \Sigma\le C\{\Pi+\sqrt{\Pi Y}\},\qquad
 \Sigma=S-32\pi^2p,\quad\Pi=P+D_2,\quad Y=(T-a^2)_+.
\end{equation}
When $a^2\ge T$, this becomes $\Sigma\le C(P+D_2)$.
Its analytic mechanism is the cancellation of a non-small compatible
background in the same actual spinor's Hessian identity, followed by
the determinant budget and a same-field weak Bianchi test. It rules
out transferring the entire cost into independent scalar relaxation.

The following exact algebraic test isolates this additional content.
It is not a metric or a monopole construction.

\begin{proposition}[Formal budgets do not force linear order]
Fix $A>0$. For $0<h\le1/4$, the formal data
\begin{align}
 p&=A,&T_+&=A(1-h),&T_-&=A(1+h)/3,\\
 P&=8\pi^2Ah^2,&S&=32\pi^2A(1+h+2h^2),\\
 W_-&=\tfrac23\pi^2Ah^2,&
 W_+&=\tfrac43\pi^2A(1-2h)+\tfrac23\pi^2Ah^2
\end{align}
satisfy both Chern--Weil identities, $W_\pm>0$, the scalar lower
bound, and $D_\beta>0$ for every $\beta>0$. Nevertheless
\begin{equation}\label{eq:formal-small}
 \J=\frac{4h^2}{1+h+3h^2}=o(h).
\end{equation}
\end{proposition}
\begin{proof}
The Chern--Weil and positivity statements follow by substitution.
Set $X=2\sqrt{1+h+2h^2}$ and
$Z=\sqrt{1-2h+h^2/2}$. Then
\[
 \frac{D_\beta}{8\pi^2A}
 =\left(\frac X{\sqrt\beta}-\sqrt\beta Z\right)^2+(X+Z)^2-9.
\]
For $0<h\le1/4$, $9-X^2-Z^2=4-2h-17h^2/2>0$, and
\[
 4X^2Z^2-(9-X^2-Z^2)^2
 =\tfrac94h^2(32-40h-25h^2)>0.
\]
Thus $2XZ>9-X^2-Z^2$, proving $(X+Z)^2>9$.
Finally $R=8\pi^2A(1+h+3h^2)$ gives \eqref{eq:formal-small}.
\end{proof}
For these data,
\begin{equation}
 D_2=108\pi^2Ah^2,\quad
 \Pi=116\pi^2Ah^2,\quad
 \Sigma=32\pi^2A(h+2h^2).
\end{equation}
Assigning the formal square $a^2=T_+$ gives $Y=0$; then
\eqref{eq:linear} fails for small $h$ with any fixed $C$.
Hence \eqref{eq:linear} is not a consequence of the \emph{listed
numerical constraints}. This does not establish realizability of the
formal data, integrality of their topology, or non-derivability from
every possible classical geometric theorem. It identifies exactly the
extra information needed from the actual equations.

In the specified actual balanced family, the resulting conclusion is
\[
 c h_j\le\alpha_j:=\inf_g\E(g)=\inf_g\J(g)
 \le\frac{4h_j}{1+h_j}.
\]
The linear scale is established in that range. It concerns
$T_+/A=1-h$, whereas the critical N1 comparison concerns $T_+/A=2/3$;
these are different parameter regions of the shared theory.

\section{The subsequent extremal construction changes the candidate value}
\label{sec:subfour}

The web comparison predates the strict finite-scale cell construction.
The 4 October continuation now supplies a once-blown-up complex-hyperbolic
ball with complete transition region included. In background
normalization $\Ric=-6g$, its explicit profile has
\begin{equation}\label{eq:cell-ball}
 \frac{J_6}{8\pi^2}:=
 \frac{1}{8\pi^2}\int|\Ric+6g|^2
 \in[0.999116243135324311\pm2.24\cdot10^{-19}]<1.
\end{equation}
The input coefficients are exact dyadic rationals. Removable bolt
singularities are eliminated algebraically before complex-ball
quadrature. Third-order matching at the outer boundary admits smooth
exact-collar replacement with an $O(\delta^3)$ change in $J_6$.
Scaling the metric by six gives the same integrated cost $J_1$ on
the $\Ric=-g$ background. The complete derivation and dated continuation
notes are retained in the accompanying \path{sources/20261004/} directory.

Here is the global implication, independent of optimizer convergence.
Let a complete cell have topology
 $C=D\#k\overline{\CP}^{\,2}\#\ell(S^1\times S^3)$, exact CH
outer collar, and $q=k+4\ell>0$. Write $d=\Vol(D)$,
$v=\Vol(C)$, $I_C=\int_Cs$, $t_C=I_C+4v$, and
$J=\int_C|\Ric+g|^2$. Inserting $n_j=o(A_j)$ copies in covers
with volume $2\pi^2A_j$ gives
\begin{gather}
 V_j=2\pi^2A_j+n_j(v-d),\\
 I_j=-4V_j+n_jt_C,\qquad R_j=4V_j+n_j(J-2t_C),\\
 \E(g_j)=\frac{4n_jJ-n_j^2t_C^2/V_j}
 {4V_j+n_j(J-2t_C)},\qquad
 \frac{\E(g_j)}{n_jq/A_j}\longrightarrow\frac{J}{2\pi^2q}.
\end{gather}
This is the inherited complete-cell transfer, including volume and
total scalar curvature, not only a local Ricci estimate.

For balanced proportions $k_j/Q_j\to2/9$ and
$\ell_j/Q_j\to7/36$, replace the blowups by the new cells and use
standard shrinking handles. A suitable cover tower supplies disjoint
fixed-radius CH balls at a positive density; $Q_j=o(A_j)$ ensures
enough sites. Thus, choosing a sufficiently small smoothing width,
\begin{equation}\label{eq:subfour}
 \limsup_j\frac{\alpha_j}{h_j}
 \le\frac{28}{9}+\frac{J_1}{9\pi^2}<3.999215<4.
\end{equation}
This disproves universal optimality of the candidate four in this
admissible family. It does not automatically apply to an arbitrary
previously fixed tower, compute the unrestricted optimum, or prove
that all near-minimizers split into these cells.

\subsection*{A further distinction in the latest continuation}
Pure-handle cells cannot achieve the analogous isolated saving. For
$C=D\#\ell(S^1\times S^3)$, the relative identities give
\begin{equation}
 J-32\pi^2\ell=
 \frac{S_C-16d}{12}+4(W_-)_C+\frac14\int_C(s+4)^2\ge0.
\end{equation}
Indeed, insertion into a CH cover and the classical scalar bound give
$S_C\ge16d$. Shrinking handles approach the threshold. However,
this does \emph{not} imply additivity of the optimal costs for mixed
blowup--handle cells: such a cell need not split along exact CH collars.
Nor is a mixed-cell infimum automatically the unrestricted global
infimum. These distinctions preserve the global variational question
without assigning a new mandatory route.

Strict curvature-energy savings through gluing have precedents, including
the Weyl-energy theorem of Malchiodi--Malizia \cite{MalchiodiMalizia}.
A Weyl saving alone does not imply \eqref{eq:cell-ball}: the scalar and
Ricci costs must be paid on the same metric. At fixed topology,
\[
 \Delta P=2\Delta(W_++W_-)+\Delta S/12.
\]
Thus the contribution to compare is the complete coupled cost and its
global transfer, not the general idea of gluing-induced energy reduction.

\section{Scope, suppliers, and the remaining global information}

Classical monopoles, special monopoles, BF nonvanishing and actual
Pin-minus lifting supply different topological and solvability inputs
\cite{Kotschick,IshidaLeBrun,Pin}. They enter common analysis only after
their quantifiers, projection budgets, class squares and covering
normalizations are satisfied. BF is not independent of monopole supply
merely because its nonvanishing theorem has a different name. A formal
twisted class is not a Pin-minus solution.

The established finite-density endpoint has the form
\begin{equation}
 p\ge A>0,\quad a^2\ge T_+,\quad\chi/A\le B<\infty
 \quad\Longrightarrow\quad
 \kappa_\beta:=\avint|\Ric+g|^2+\frac{D_\beta}{V}
 \ge c(B)\frac{(\beta-2)^2}{\beta^2},
\end{equation}
under its original actual conformal-supply assumptions. The manuscript
also gives a nonoptimal polynomial lower control on $c(B)$.
This is different from the single-solution finite-parameter estimate,
which uses $D_\beta/(\beta V)$. Every fixed finite density is covered;
a density-independent constant still requires additional information.
In particular the bounded-density balanced family is already within
the former theorem's scope.

\begin{center}\small
\begin{tabularx}{\textwidth}{@{}p{0.29\textwidth}X@{}}\toprule
Statement & Mathematical position after the 4 October continuation\\\midrule
Subcritical normalized bound \eqref{eq:old-gap} &
Direct consequence of classical joint curvature estimates.\\
Recovery upper bound \eqref{eq:old-upper} &
Inherited classical construction, not a newly supplied lower bound.\\
N1 and square-loss versions &
Uniform quantitative non-saturation from the stated actual-solution theorem.\\
I40/I41 linear coercivity &
Additional same-field constraint, not implied by the displayed classical numerical budgets.\\
Linear order in the balanced range &
Established; the best first coefficient is not determined.\\
Universal candidate coefficient four &
Disproved by the finite-scale construction on suitable balanced cover towers.\\
O10 &
Every fixed finite-density version established; density-independent version open.\\
Full BF classical comparison &
Known supplied-class comparison is available; the complete class budget on the different BF families remains a separate issue.\\
Exact realization &
Existing conditional realization theorems remain valid; vanishing gap alone is not an existence theorem.\\\bottomrule
\end{tabularx}
\end{center}

The common scientific statement is therefore precise: where classical
budgets degenerate or fail to determine the defect scale, the actual-field
equations can provide further uniform coercivity; real constructions
test its scale, sharpness and geometric scope. The single-cell result
is one such test, not a replacement for the all-metric theory.

This comparison supplements ongoing research; it does not set an order
of tasks or prescribe a decomposition, symmetry class, numerical ansatz
or proof technology. A new identity, a broader actual supplier, a sharp
global bound, a counterexample or an exact realization may each change
the mathematics. Local information becomes a global conclusion only
through a proved implication at the required scale.

\section*{Reading connections}
Paths below are relative to the Geometric Analysis Skill root.
\begin{itemize}
\item Conventions and current statements:
\path{references/conventions-and-literature.md} and
\path{references/results-and-questions.md}.
\item N1: \path{research/paper/sections/tracefree_gap.tex}.
Linear analysis: \path{research/paper/sections/linear_coercivity.tex}
and \path{research/paper/sections/single_solution_coercivity.tex}.
\item Finite-density endpoint:
\path{research/paper/sections/finite_density_endpoint.tex} and
\path{research/paper/sections/web_primitive_rate.tex}.
\item Global extremal reduction:
\path{research/paper/sections/variational_geometry.tex}.
Complete-cell transfer:
\path{research/sources/continuation/20260928_relative_cell_cost.md},
Sections 1--2.
\item Latest construction and its further distinctions:
\path{research/supplements/sources/20261004/}, containing the original
construction and review texts, frozen coefficients and integration code.
\end{itemize}
This supplement derives the comparison formulas and formal-budget test
above. N1, I40/I41 and the endpoint are used with their original proofs
and inputs; they are not re-proved by a literature comparison. The latest
construction is a dated addition, not a claim to have exhausted related
literature or all subsequent research.

\begin{thebibliography}{9}
\bibitem{LeBrunRicci} C. LeBrun,
\emph{Ricci Curvature, Minimal Volumes, and Seiberg--Witten Theory},
\href{https://arxiv.org/abs/math/0003068}{arXiv:math/0003068},
especially Sections 2--4 and Theorem 4.3.
\bibitem{LeBrunConvex} C. LeBrun,
\emph{Four-Manifolds, Curvature Bounds, and Convex Geometry},
\href{https://arxiv.org/abs/math/0611450}{arXiv:math/0611450} (2006),
Theorem A and Theorem 4.10.
\bibitem{LeBrunOverview} C. LeBrun,
\emph{An Overview of Einstein 4-Manifolds},
\href{https://arxiv.org/abs/2609.03033v1}{arXiv:2609.03033v1} (2026),
Theorems 7.4--7.5.
\bibitem{LeBrunOptimal} C. LeBrun,
\emph{Curvature Functionals, Optimal Metrics, and the Differential
Topology of 4-Manifolds},
\href{https://arxiv.org/abs/math/0404251}{arXiv:math/0404251}, Section 11.
\bibitem{IshidaLeBrun} M. Ishida and C. LeBrun,
\emph{Curvature, Connected Sums, and Seiberg--Witten Theory},
\href{https://arxiv.org/abs/math/0111228v2}{arXiv:math/0111228v2},
especially Section 5.
\bibitem{Kotschick} D. Kotschick,
\emph{Monopole Classes and Einstein Metrics},
\href{https://arxiv.org/abs/math/0306022}{arXiv:math/0306022}.
\bibitem{Pin} M. Ishida, S. Matsuo and N. Nakamura,
\emph{Yamabe Invariants and the $\mathrm{Pin}^{-}(2)$-Monopole Equations},
\href{https://arxiv.org/abs/1505.05948}{arXiv:1505.05948}.
\bibitem{MalchiodiMalizia} A. Malchiodi and F. Malizia,
\emph{Weyl Energy and Connected Sums of Four-Manifolds},
\href{https://arxiv.org/abs/2505.17752v1}{arXiv:2505.17752v1},
Theorem 1.1.
\end{thebibliography}
\end{document}
